\section*{ID8484: TOROIDAL KNOT TOPOLOGY Canonical Statutory Formula: K\_DNA Cong T(3, 2) ⊗ S\textasciicircum{}1}
\paragraph{Metadata.}
PiID: 3441375223970 | RTEID: RTE:P32945.2561:T4.0182 | UUID: 2f3ffcac-d8f3-5789-9e9f-8f586fdb68c6

\paragraph{Canonical Statement.}
[ID 1176] TOROIDAL KNOT TOPOLOGY Canonical Statutory Formula: \$K\_\{DNA\} \textbackslash{}cong T(3, 2) \textbackslash{}otimes S\textasciicircum{}1\$ Formal Technical Dissertation: Defines the topological manifold of the operator. The "DNA Helix Operator" is geometrically modeled not as a simple cylinder, but as a \$(3,2)\$ Torus Knot tensor-producted with a circle (\$S\textasciicircum{}1\$), creating a higher-dimensional braiding structure essential for trapping wavefunctions. Source: [2]

\paragraph{Canonical Equation.}
\[
K_{DNA} \cong T(3, 2) \otimes S^1
\]

\paragraph{Dictionary.}
TOROIDAL KNOT TOPOLOGY Canonical Statutory Formula: K\_DNA cong T(3, 2) ⊗ S\textasciicircum{}1 Formal Technical Dissertation: Defines the topological manifold of the operator.

\paragraph{Encyclopedia.}
TOROIDAL KNOT TOPOLOGY Canonical Statutory Formula: K\_DNA Cong T(3, 2) ⊗ S\textasciicircum{}1 is a symbolic expression, written K\_DNA cong T(3, 2) ⊗ S\textasciicircum{}1. It introduces a symbol or relation used elsewhere in the framework. It is built from K\_DNA, T, S. Within the Trawin Zero Point registry it serves as one element of the framework's mathematical narrative.
